\begin{recipe}{Pesto genovese}
\label{recipe:pesto-genovese}
\kicker[Bijgerecht,Vegetarisch,Italiaans,Basisrecept]{Bijgerecht · vegetarisch · Italiaans · basisrecept}
\index[register]{Pesto genovese}
\index[register]{Basilicum}
\index[register]{Pijnboompitten}
\index[register]{Parmezaanse kaas}
\index[register]{Pecorino}
\index[register]{Knoflook}

\meta{45 minuten \dvd Circa 250 g pesto}

\marginimage{images/pesto-genovese}
\lettrine{D}{e} potjes pesto uit het schap in de supermarkt hebben weinig te
maken met verse pesto: die smaken vooral naar azijn en verduurzaamde olie. De
verse pesto uit de koeling is al een stuk beter, maar zelf maken smaakt toch
het lekkerst, dankzij verse basilicum en olijfolie die niet wekenlang hoeft
te worden bewaard.

\blockrule{Ingrediënten}
\begin{ingredients}
  \ing{4 handenvol}{verse basilicumblaadjes}\index[register]{Basilicum}
  \ing{30 g}{pijnboompitten}\index[register]{Pijnboompitten}
  \ing{50 g}{oude Parmezaanse kaas, geraspt}\index[register]{Parmezaanse kaas}
  \ing{30 g}{Pecorino, geraspt}\index[register]{Pecorino}
  \ing{1 teentje}{knoflook, gepeld}\index[register]{Knoflook}
  \ingb{grof zeezout, naar smaak}
  \ingb{extra vierge olijfolie, naar behoefte}
\end{ingredients}

\blockrule{Bereiding}
\tip{Traditioneel wordt pesto met een vijzel gestampt in plaats van
gepureerd, dat geeft een grovere, romigere textuur. Wij hebben geen vijzel in
huis, dus gaat het bij ons met de staafmixer: pulseer kort en vaak in plaats
van continu, anders wordt hij heet en verliest de basilicum zijn frisse
groene kleur.}
\begin{steps}
  \item Was de basilicumblaadjes en dep ze goed droog met keukenpapier.
  \item Rooster de pijnboompitten zonder vet in een steelpannetje op
        middelhoog vuur, 2 tot 3 minuten tot ze goudbruin en geurig zijn, als
        je ze geroosterd wilt.
  \item Doe de basilicum, pijnboompitten, knoflook, een snuf grof zeezout en
        een scheut olijfolie in een hoge, smalle kom.
  \item Pulseer met de staafmixer tot een grove pasta en voeg geleidelijk
        meer olijfolie toe, tot de pesto glad en romig is.
  \item Roer de Parmezaanse kaas en de Pecorino er met een vork doorheen: dat
        geeft een mooiere textuur dan alles in één keer met de staafmixer
        pureren.
\end{steps}
\end{recipe}
